\documentclass[11pt]{article}
% Degrees of separation -- main text for EPJ Data Science.
% Every number below is copied from a result write-up of a seeded script in the repository (scripts 01--64; see SI
% and paper/CLAIMS_LEDGER.md), from the interim tables those scripts write, or from code run on those tables:
% scripts/70_paper_figures.py (figures; with --si-checks, the numbers marked "SI check").
\usepackage[margin=1in]{geometry}
\usepackage[T1]{fontenc}
\usepackage{lmodern}
\usepackage{amsmath,amssymb}
\usepackage{graphicx}
\usepackage{booktabs}
\usepackage{array}
\usepackage{xcolor}
\usepackage[round]{natbib}
\usepackage{hyperref}
\hypersetup{colorlinks=true,linkcolor=blue!50!black,citecolor=blue!50!black,urlcolor=blue!50!black}

% Figures are drawn at the printed text width (165 mm) and included at 1:1; the file compiles if one is missing.
\newcommand{\paperfig}[2][\textwidth]{%
  \IfFileExists{../outputs/figures/#2.pdf}{\includegraphics[width=#1]{../outputs/figures/#2}}{%
  \IfFileExists{../outputs/figures/#2.png}{\includegraphics[width=#1]{../outputs/figures/#2}}{%
  \fbox{\parbox{0.9\linewidth}{\centering\small Figure file \texttt{\detokenize{#2}} not found.}}}}}
\renewcommand{\topfraction}{0.9}
\renewcommand{\bottomfraction}{0.8}
\renewcommand{\textfraction}{0.07}
\renewcommand{\floatpagefraction}{0.75}
\setcounter{topnumber}{3}
\setcounter{totalnumber}{4}
\setlength{\emergencystretch}{3em}
\newcommand{\SI}[1]{SI Section~S#1}
\newcommand{\ci}[2]{$[#1,\allowbreak\,#2]$}
\newcommand{\cic}[2]{$\langle #1,\,#2\rangle$}

\title{Degrees of separation: graduate earnings track institutional status more than the academy's departmental hierarchy}
\author{Sora Yng}
\date{}

\begin{document}
\maketitle

\begin{abstract}
The academy ranks its own departments through faculty hiring: which departments place their PhDs as
faculty at which others. We ask whether this hierarchy carries information about the earnings of
bachelor's graduates, field by field. Using public data only, we measure within-field coupling, the
rank correlation across institutions between a department's position in the US faculty-hiring
hierarchy and its graduates' median earnings, for 52 US fields (College Scorecard, with Census PSEO as a
second, overlapping source) and, with a UK hiring hierarchy built from ORCID-derived records, for 33 UK
subjects (DfE LEO). The analysis is descriptive. Coupling is positive and varies across fields, but most of
it is shared with institution-level status. Controls for selectivity, Pell share, control and state lower
mean US coupling from $+0.43$ to $+0.14$. Within institutions, once each field's earnings slopes on
selectivity and academia-wide prestige are allowed for, one standard deviation of department prestige goes
with about 1\% higher pay (0.84\%, 95\% interval 0.00--1.64\%, for four-year earnings; up to about 3\% after
correcting for measurement error at the lower reliability bound), and field prestige predicts pay no better
than academia-wide prestige. In the UK, 66--90\% of coupling survives a direct control for graduates' prior
attainment, and what survives tracks intake selectivity rather than hiring-network prestige. Within fixed
graduation cohorts, coupling rises by $+0.031$ per year in the US and $+0.026$ in the UK, a change that mixes
career and calendar time (in the US the career-time part lies between $+0.014$ and $+0.032$ per year if
calendar and cohort effects are non-negative); in a weakly identified US decomposition the rise loads on
academia-wide rather than field prestige. Coupling is weak in UK subjects paid on national scales (nursing,
medicine, teaching: $+0.12$ against $+0.52$), a pattern known from UK returns, and across US fields the
prestige--pay association net of geography falls with the share of graduates in strictly licensed
occupations, a contrast between discipline families that public data cannot separate from field type or
employer sector. The academy's departmental hierarchy adds little labour-market information beyond
institution-wide status.
\end{abstract}

\noindent\textbf{Keywords:} faculty hiring networks; academic prestige; graduate earnings; college
selectivity; field of study; administrative data; career dynamics; science of science

% =====================================================================================
\section{Introduction}\label{sec:intro}

The academy ranks its own departments. Each time a department hires a professor it reveals which
doctoral programmes it trusts to train faculty, and aggregated over a decade these choices form a steep
hierarchy in which a few institutions train most of the faculty in their field
\citep{clauset2015systematic, wapman2022quantifying}. The resulting prestige ordering is estimated from
peers' hiring decisions, not from admissions statistics, reputation surveys or composite rankings, and
it exists field by field across the US academy. We ask whether this hierarchy carries information about how
the academy's bachelor's graduates fare outside it. Within a field, do graduates of departments higher
in the hiring hierarchy earn more? If they do, does the association belong to the department or to the
institution in which the department sits?

The question matters for two literatures. For the science of science, the faculty-hiring hierarchy is
a carefully measured status ordering, built from a census of hires, and position in it shapes productivity, the
spread of ideas and who becomes faculty \citep{way2019productivity, morgan2018prestige,
morgan2022socioeconomic, zhang2022labor}. Whether the ordering is also legible outside academia, where
nearly all bachelor's graduates work, has not, to our knowledge, been measured. For the economics and sociology of higher
education, the question is a sharper version of an old one. Earnings rise with institutional
selectivity and brand, and much of that gradient is attributed to who enrols
\citep{dale2002estimating, dale2014estimating, mountjoy2021returns}. The size of the status gradient in pay
differs by field of study \citep{eide2016where, quadlin2021same, hastings2013some, bloem2024understanding}, and in
the United Kingdom an institution's academic standing adds little information about graduates' returns
beyond the selectivity of its intake \citep{hussain2009university, britton2022how}. In employer-learning
models the weight employers place on an easily observed credential falls as they learn about
productivity, relative to hard-to-observe correlates of productivity \citep{farber1996learning,
altonji2001employer}, while studies of elite hiring find employers screening on pedigree
\citep{rivera2015pedigree}. A department-level, peer-defined prestige axis lets us ask how much of the
status gradient in pay belongs to the department, and how the gradient changes over careers.

We measure within-field coupling: the Spearman rank correlation, across the institutions that teach a
field, between a department's position in the US faculty-hiring hierarchy \citep{wapman2022quantifying}
and the median earnings of its bachelor's graduates. We map coupling across 52 US fields on College
Scorecard earnings \citep{scorecard} and repeat it on Census Post-Secondary Employment Outcomes (PSEO)
\citep{pseo}, a second earnings source that covers many of the same institutions. For the United Kingdom
we build an academia-wide hiring hierarchy from ORCID-derived records of PhD-to-faculty moves
\citep{orcid} and set it against the Department for Education's Longitudinal Education Outcomes (LEO)
earnings for 33 subjects, which are also released by graduates' prior attainment. All inputs are
public, and every result is descriptive.

Four findings hold. First, coupling is positive, varies across fields, and has a similar field ordering on
the two US earnings sources. Second, most of it is shared with institution-level status. Controls for
institutional selectivity, Pell share, control and state lower mean US coupling from $+0.43$ to $+0.14$;
within institutions, one within-field standard deviation (SD) of department prestige goes with about 1\%
higher four-year earnings (0.84\%, 95\% interval 0.00--1.64\%) once each field's slopes on selectivity and
academia-wide prestige are allowed for; and a field's own prestige ordering predicts pay no better than the
academia-wide ordering. Third, in the UK 66--90\% of coupling survives a direct control for graduates' own
prior attainment, and what survives tracks intake selectivity rather than hiring-network prestige. Fourth,
within fixed graduation cohorts coupling rises, by $+0.031$ per year in the US and $+0.026$ per year in the
UK. How much of the US rise is career time rather than calendar time can only be bounded, part of it may
reflect the timing of graduate training, and in a weakly identified decomposition it loads on the
academia-wide ordering, which tracks selectivity, rather than on field prestige. At the other end, the
prestige--pay association is weak in UK subjects on national pay scales and falls across US fields with the
share of graduates in strictly licensed occupations, contrasts between discipline families that public data
cannot separate from field type or employer sector.

Much of this is known in kind. That the status gradient in pay differs across fields is documented by
\citet{bloem2024understanding} on an earlier College Scorecard release and by \citet{britton2022how} on UK
administrative data; that academic quality adds little beyond selectivity is reported for the UK by
\citet{britton2022how} and \citet{hussain2009university}; and the direction of the career-time result is
known from descriptive work \citep{thomas2005postbaccalaureate, macleod2017big}, although causal designs
find small or fading premia \citep{bordon2020employer, mountjoy2021returns}. Our cross-field map and our UK
attainment check reproduce these patterns with a different prestige axis. What is new is (i) the
faculty-hiring hierarchy as the prestige axis, which separates a department's standing in its field (field
prestige, $F$) from its institution's standing in the academy as a whole (academia-wide prestige, $G$);
(ii) with it, a within-institution test of whether departments of higher standing inside the same
institution have better-paid graduates, which is the test the question of the title requires; (iii)
fixed-cohort estimates, in two countries, of how coupling changes with years since graduation; and (iv) two
measurement contributions for data science: a bracket on the sampling reliability of the published
SpringRank ranks, built from the suppressed public edge list, and an academia-wide UK hiring hierarchy built
from ORCID-derived records.

For the academy's departmental hierarchy, the answer is mostly negative. The hierarchy tracks graduate
pay, but largely through the institutional status it shares with selectivity. Beyond institution-level
status the department ordering goes with a small remainder, about 1\% of pay per SD within institutions,
which is individually clear only in computer science; whether it varies across fields is fragile, and
public data cannot attribute it to department standing or to program-level selection.
Section~\ref{sec:background} positions the study, Section~\ref{sec:data} describes data and measurement,
Section~\ref{sec:results} reports the results, Sections~\ref{sec:discussion} and~\ref{sec:limits} discuss
them and their limits, and Section~\ref{sec:methods} states the main estimators. Further methods are in the
Supplementary Information (SI).

% =====================================================================================
\section{Background and positioning}\label{sec:background}

\paragraph{Returns to college quality and selection.} Once application behaviour absorbs selection, the
average earnings payoff to a more selective college is small, with larger returns for disadvantaged
students \citep{dale2002estimating, dale2014estimating}. Regression-discontinuity designs find sizeable
effects in some settings \citep{hoekstra2009effect}, and selection-corrected models still find large
payoffs to highly selective colleges \citep{brewer1999does, witteveen2017earnings}. Single quality
proxies are noisy and understate the return \citetext{\citealp{black2006estimating}; see also \citealp{black2004robust}},
returns differ across the earnings distribution \citep{andrews2016quantile}, raw earnings differences between Texas
colleges shrink sharply once selection is controlled \citep{cunha2014measuring}, and selectivity predicts
value added poorly within students' choice sets \citep{mountjoy2021returns}. The payoff to elite
colleges is concentrated in the upper tail of outcomes \citep{zimmerman2019elite, chetty2023diversifying},
where elite social networks inside colleges also matter \citep{michelman2022old}. Most of this work
measures institutional standing by selectivity, rankings, brand or value added, in earnings or in learning
\citep{riehl2019learning}. The closest use of an academic measure is \citet{hussain2009university}: in UK
graduate surveys an institution's research rating predicts wages (about 4\% per SD with controls), but it
correlates 0.87 with entry tariff and turns negative, imprecisely, when entered alongside other quality
measures.

\paragraph{Field of study and status by field.} Returns to field of study are large and heterogeneous.
In Norway, admission cutoffs reveal widely different payoffs across fields net of institution quality
\citep{kirkeboen2016field}, and US studies document large lifetime and trajectory differences by major
\citep{altonji2012heterogeneity, altonji2016analysis, webber2014lifetime, kim2015field,
andrews2024returns}. The size of the institutional-status gradient in pay also varies by field. Earnings gaps across
selectivity levels are largest for business and smallest for science majors \citep{eide2016where};
after matching on US survey data, only business and social-science majors show a selectivity premium
\citep{quadlin2021same}; and Chilean returns to admission are large for highly selective degrees and for
health, science and social-science degrees \citep{hastings2013some}. Closest to us in the US,
\citet{bloem2024understanding} show descriptively on an earlier release of College Scorecard field-of-study
data, our earnings source, that median pay barely varies with admit rate in nursing, teacher education,
social work and criminal justice but varies strongly with it in computer science, finance, economics,
political science and marketing; cross-classified models on Scorecard earnings find selectivity matters
more outside STEM \citep{muse2024college}. Our data reproduce their ordering of these fields, though not
the flatness: the within-field rank correlation of selectivity (minus the admit rate) with four-year
earnings is $+0.62$ to $+0.77$ in their five steep fields and $+0.25$ to $+0.37$ in the three flat fields
we cover (\texttt{scripts/55}). In the UK, returns by institution and subject were earlier estimated on
survey data, where much of the variation across courses is attributed to selection
\citep{walker2011differences, walker2018university}, and are now estimated on administrative LEO records
\citep{belfield2018impact, britton2019improving, britton2020impact}. The closest design to ours is
\citet{britton2022how}. Within subject, returns at age 30 track the prior attainment of a degree's intake
closely in law, business and economics (correlations 0.93, 0.89 and 0.79) and not at all in medicine (0.09)
or education ($-0.08$), which the authors attribute to careers with centrally regulated wages; net of
selectivity, no subject's league-table rating, which includes research intensity, correlates positively and
significantly with returns.

\paragraph{Faculty-hiring hierarchies.} In sociology, centrality in the PhD-exchange network reproduces
84\% of the variance in departments' reputational prestige \citep{burris2004academic}. Across fields,
hiring networks are steeply hierarchical, with a few institutions training most faculty
\citep{clauset2015systematic, wapman2022quantifying}. PhD origin is tightly linked to faculty placement
\citep{fowler2007social, headworth2016credential}, placement records have been used to rank departments
\citep{amir2008ranking}, and hiring networks divide into clusters with little hiring between them
\citep{tervio2011divisions}. Position in the hierarchy shapes productivity, the diffusion of ideas and who
becomes faculty \citep{way2019productivity, morgan2018prestige, morgan2022socioeconomic, zhang2022labor}.
The hierarchy is estimated with SpringRank, which orders institutions so that most hires flow down the
ordering \citep{debacco2018physical}. This literature measures standing within academia. To our knowledge it
has not asked whether position in the hierarchy predicts the pay of the academy's bachelor's graduates in
the wider labour market, a question made less than obvious by the weak link between research and teaching
quality \citep{hattie1996relationship}.

\paragraph{Employer learning and career dynamics.} In employer-learning models the wage weight on an easily
observed credential falls with experience relative to hard-to-observe correlates of productivity
\citep{altonji2001employer}, while the unconditional association of schooling with the level of wages need
not change \citep{farber1996learning}. Employers learn quickly \citep{lange2007speed, arcidiacono2010beyond},
learning varies by occupation \citep{mansour2012does}, and pay dispersion grows with experience mostly
because productivity diverges \citep{kahn2014employer}. Prestige premia fade with experience in Chile
\citep{bordon2020employer}, and university credentials matter less for later promotions in Japan
\citep{araki2016university}. Descriptive evidence points the other way: US graduates of higher-quality
colleges have faster earnings growth in their first four years \citep{thomas2005postbaccalaureate}, and
college reputation correlates with earnings growth in Colombia \citep{macleod2017big}; selective-college
payoffs are large at both four and ten years after graduation in two US surveys
\citep{witteveen2017earnings}. Field-specific and entry-condition forces operate as well
\citep{roksa2010what, oreopoulos2012short, altonji2016cashier, deming2020earnings}. Separating career time
from calendar time and cohort requires assumptions \citep{fosse2019analyzing, fosse2019bounding}, and
administrative earnings records lose graduates selectively \citep{foote2026attrition}.

\paragraph{What is known and what is new.} Known: small average returns to selectivity net of selection,
elite effects in the tails, fast employer learning, and returns to status that vary by field in the US,
Chile and the UK, with little pay gradient in UK medicine and education, where wages are set centrally
\citep{britton2022how}, and in US nursing, teacher education and social work \citep{bloem2024understanding}.
Our cross-field map and our UK check, which uses a coarser, band-level attainment control than
\citet{britton2022how}, are replications in kind; what is new is listed at the end of
Section~\ref{sec:intro}. Licensing, pay schedules and employer sector enter only in
Section~\ref{sec:boundary}, as candidate readings of where coupling is weak.

% =====================================================================================
\section{Data and measurement}\label{sec:data}

Table~\ref{tab:data} lists every source with its unit, the sample used and its known limitations. All
inputs are public; versions, URLs and access dates are recorded in the repository (\SI{1}).

\subsection{Prestige: the faculty-hiring hierarchy}\label{sec:prestige}

We use the published prestige ranks of \citet{wapman2022quantifying}, estimated with SpringRank from a
census of faculty at 368 PhD-granting US institutions, 2011--2020 \citep{wapman2022data} (the published
academia-wide ranks have 371 rows, 369 of them named). For institution $i$ and field $f$, field
(department) prestige $F_{if}$ is minus the published field rank, and academia-wide prestige $G_i$ is minus
the published rank in the network that pools hires across all fields. Both are fixed across earnings
horizons, so any change over careers is on the earnings side. $G$ is sometimes read as institutional brand;
we use that word only as an interpretation, because $G$ also tracks selectivity. Across 242 US universities
its Spearman correlation with average SAT is $+0.68$, and across 133 UK providers the UK analogue correlates
$+0.71$ with intake selectivity. $F$ and $G$ are themselves close: within the field-by-cohort cells of the
PSEO panel their Spearman correlation averages 0.83 (range 0.55--0.95).

The public release omits institution pairs linked by a single hire (the smallest edge weight in the
public edge list is two). In the 40 fields with published hire counts, the public field-level edges
carry a median 15\% of the field's US-doctorate faculty (self-hires excluded), against 67\% for the
academia-wide network, so a SpringRank rebuilt from public edges reproduces the published field ranks
only partly (median Spearman 0.73, range 0.32--0.87; academia-wide 0.908). We therefore use the
published ranks, whose sampling reliability public data can only bracket (Section~\ref{sec:methods}).
Measured on public edges, the split-half reliability of $F$ averages 0.81 across 57 fields (range
0.56--0.95), a lower bound; extrapolated to the number of hires behind the published ranks it is 0.96. The
reliability of $G$ is 0.99 on the fields' samples (0.96 across the whole academia-wide network; \SI{2}).
Where reliability matters, results are reported at both ends of the bracket.

\subsection{Earnings and placement}\label{sec:earnings}

\paragraph{US: College Scorecard.} The canonical outcome is the median earnings of Title IV--aided
bachelor's completers four years after completion, from the Field-of-Study file released in June 2026
\citep{scorecard}, aggregated over each field's four-digit CIP codes through a curated map of 66 fields
(\SI{1}). Matching to the published ranks gives 3,444 institution-by-field cells in the 52 fields with
at least 15 matched institutions (median 48 institutions per field, range 15--173; \texttt{scripts/55}), the canonical
sample, and 3,500 cells in the 57 fields with at least 8, which we use for the heterogeneity statistics, the
reliability bracket, the comparison of $F$ with $G$ and the cross-source comparison. Scorecard covers the
ranked institutions about evenly across academia-wide prestige quartiles (74.2\%, 65.6\%, 71.7\% and
65.6\% on the 66-field map, top to bottom; \SI{1}). Institutional covariates come from the Scorecard
``most recent'' institution file (2026 build): average SAT, admission rate, Pell share, control,
institution-wide median earnings ten years after entry, and the earnings level of the institution's state.
These describe recent entering classes, not the cohorts whose earnings we observe
(Section~\ref{sec:limits}). No public US file measures the entry attainment of a program's own students;
the Field-of-Study file has 178 columns and none carries SAT or ACT information. One-year and five-year
medians from the same release refer to other graduating cohorts and are used only as replications.

\paragraph{US: PSEO.} PSEO reports unemployment-insurance wage earnings of graduates, in constant 2023
dollars, one, five and ten years after graduation (y1, y5, y10) by three-year graduation cohort
\citep{pseo}. We use release V4.13.0 (2025Q4; 952 institutions) as canonical and the August 2026 release
V4.14.1 (2026Q2; 1,117 institutions) as a re-check. For career time we build fixed cohorts (graduates of
2001--03, 2004--06, 2007--09 and 2010--12) and keep, within each field and cohort, the institutions
observed at all three horizons, with at least 15 per cell: 106 institutions, 142 field-by-cohort cells
and 39 fields on V4.13.0, and 116, 145 and 40 on V4.14.1. PSEO partners are mostly public institutions:
only 20 of the 93 top-quartile ranked institutions have a usable five-year cell on V4.13.0 (26 on
V4.14.1), so the prestige range is truncated at the top. PSEO Flows give the employed graduates of each
institution-by-CIP-2 cell by employer sector. We define a market share $M$ (information, finance and
professional services; NAICS 51, 52, 54) and a setting share $S$ (education, health care and public
administration; NAICS 61, 62, 92), and a family prestige $P$, the completions-weighted mean of within-field
prestige percentiles over the ranked fields of each CIP-2 family.

\paragraph{UK: LEO and an academia-wide hiring hierarchy.} The DfE provider-level LEO file (``Graduate
outcomes (LEO): provider level data'', 2022-23 release) reports median earnings by provider, CAH2 subject,
graduating cohort and year after graduation (YAG 1, 3, 5), and separately by prior-attainment (UCAS tariff)
band. We use the cohorts 2013/14 to 2016/17 and the 33 subjects with at least 15 providers. For prestige
we run SpringRank on 11,987 PhD-to-faculty moves between institutions in Great Britain in the ORCID-derived
academic mobility edges (Zenodo, doi:10.5281/zenodo.19651302, version 20260419), which are built from ORCID
public records \citep{orcid} with organisations resolved to ROR identifiers; institutions are matched to
providers through the ROR data dump v2.8 (Zenodo record 20512981). The network has 551 institutions; 149
match LEO providers and 134 form the primary sample (network degree at least 5; \SI{1}). The resulting
academia-wide ordering places Russell Group members at a mean percentile of 0.86 against 0.43 for other
providers (AUC 0.93); its split-half correlation is 0.85 for half-networks and its reliability 0.92 for the
full network (Spearman--Brown). These are internal checks. Its external validity is less certain: in the
US, an academia-wide ranking rebuilt from ORCID records agreed with the census-based published ranks at
Spearman 0.74, against 0.91 for SpringRank on the public census edges (legacy \texttt{scripts/10}--11;
\SI{2}). Field-tagged UK networks are too sparse (47 to 504 edges per usable field) to serve as the main
axis, so the UK analysis has no department axis. Intake selectivity is the institution's share of graduates
who entered with at least 360 tariff points.

\paragraph{Other sources.} Opportunity Insights college-level outcomes \citep{chetty2020income} give,
for 242 PhD-granting universities linked to single-college units, median earnings and the share of
former students in the top 1\% of their cohort's earnings, with parental income and selectivity. The
2023 American Community Survey 1-year public-use microdata (U.S. Census Bureau) give field-level shares of
degree holders in licensed occupations and in government or non-profit employment.

\subsection{Coupling and inference}\label{sec:coupling}

Coupling is $\rho_f=\mathrm{Spearman}_i(F_{if},Y_{if})$ across the institutions $i$ that teach field
$f$, with $Y$ median earnings. As a rank statistic within field it ignores pay differences between fields
and asks only whether the prestige and pay orderings coincide; it does not measure how large the pay
differences between high- and low-prestige institutions are. Earlier drafts reported the gap $1-\rho_f$,
which carries the same information. Partial couplings correlate the rank residuals of $F$ and $Y$ after
regressing both on the ranks of the controls. A field is flagged as reliable when its earnings carry
cross-institution signal beyond cohort sampling noise (signal fraction at least 0.5 under a noise model
calibrated to cohort sizes), when the bootstrap interval of $\rho_f$ does not overlap the band that
sampling noise alone would produce under perfect agreement, and when it has at least 10 institutions; 20 of
the 52 canonical fields pass. We use the flag as a check, not a filter, because the signal fraction is
miscalibrated for some fields (\SI{3}).

Fields share institutions, so cross-field statistics are not independent, and each interval answers a
stated question. Means across fields on Scorecard use a joint institution bootstrap that resamples
institutions once per draw for every field; these intervals are conditional on the observed fields. PSEO
career-time statistics are about fields in general and use the two-way (field, institution) cluster
variance, which treats both fields and institutions as sampled; beside it we report the crossed
field-by-institution bootstrap, which is conservative for these statistics (under a within-cell null its
variance is 3.0--3.7 times the true sampling variance). UK intervals resample subjects and then providers
within subject. Field-level moderators are tested by field-label permutation, which treats fields as
exchangeable, and also with a bootstrap that resamples whole CIP-2 families, shown in angle brackets.
Brackets are 95\% intervals throughout. The canonical specification was fixed before the manuscript was
drafted, after the analyses had been run; alternatives are collected in a specification curve
\citep{simonsohn2020specification} in \SI{4}, and \SI{9} lists every hypothesis the analyses examined. Every
number comes from a seeded script in the repository or, where marked ``SI check'', from code run on the
repository's tables; the SI maps each number to its source.

\begin{table}[!tbp]
\centering
\caption{Data sources and samples. All sources are public; provenance is recorded in
\texttt{data/raw/SOURCES.md} of the repository.}\label{tab:data}
\footnotesize
\renewcommand{\arraystretch}{0.95}
\begin{tabular}{@{}>{\raggedright\arraybackslash}p{2.4cm}>{\raggedright\arraybackslash}p{3.3cm}>{\raggedright\arraybackslash}p{3.3cm}>{\raggedright\arraybackslash}p{2.2cm}>{\raggedright\arraybackslash}p{3.5cm}@{}}
\toprule
Source & Unit and release & Sample used here & Role & Known limitation \\
\midrule
Faculty-hiring ranks \citep{wapman2022data} & institution $\times$ field; published field and academia-wide SpringRank ranks, 2011--2020 & 66 mapped fields; 57 with $\ge 8$ and 52 with $\ge 15$ matched institutions & $F$, $G$ & public edges omit single-hire pairs; reliability of $F$ bracketed (0.81--0.96) \\
College Scorecard, Field of Study \citep{scorecard} & institution $\times$ CIP-4, bachelor's; median earnings 4 years after completion (1, 5 years as replications); June 2026 & 3,444 cells, 52 fields & $Y$ (canonical) & Title IV completers; medians only; no program-level entry attainment \\
College Scorecard, institution file \citep{scorecard} & institution; ``most recent'' file, 2026 build & 46--47 fields with SAT data; 52 with admission rate & selectivity, Pell share, control, institution earnings, state level & institution-level; describes recent entering classes, not the earnings cohorts \\
PSEO earnings, V4.13.0 and V4.14.1 \citep{pseo} & institution $\times$ CIP-4 $\times$ 3-year cohort; p25, p50, p75 at y1, y5, y10 & fixed cohorts 2001--2012: 106 (116) institutions, 39 (40) fields & second source; career time & public-skewed; 20 (26) of 93 top-quartile institutions; UI-covered jobs \\
PSEO Flows \citep{pseo} & institution $\times$ CIP-2 $\times$ employer sector & 1,312 cells at 108 ranked institutions; 4,773 cells at 348 institutions with SAT; 21 families & placement shares $M$, $S$ & noise-protected counts; employer industry, not occupation \\
DfE LEO provider file, 2022-23 release & provider $\times$ CAH2 $\times$ cohort $\times$ YAG; also by prior-attainment band & 33 subjects, 134 providers, cohorts 2013/14--2016/17, YAG 1/3/5 & UK earnings & medians; small cells suppressed; coarse bands \\
ORCID-derived mobility edges (Zenodo, doi:\allowbreak\,10.5281/\allowbreak zenodo.\allowbreak 19651302), from ORCID \citep{orcid}; ROR v2.8 & PhD-to-faculty moves within Great Britain & 11,987 moves, 551 institutions & UK $G$ & registry self-report; academia-wide only \\
Opportunity Insights college tables \citep{chetty2020income} & college; outcomes by parent income & 242 PhD-granting universities & tail outcomes & pools all fields of a college \\
ACS PUMS 2023, 1-year (U.S. Census Bureau) & person & field-level shares & licensed and public or non-profit employment shares & shares include graduate-degree holders \\
\bottomrule
\end{tabular}
\end{table}

% =====================================================================================
\section{Results}\label{sec:results}

We report the four findings in turn and then the fields where coupling is weak.
Table~\ref{tab:ladder} collects the selectivity controls for both countries and
Table~\ref{tab:career} the career-time estimates.

\subsection{Coupling is positive, heterogeneous and similar on two earnings sources}\label{sec:map}

Graduate pay tracks the hiring hierarchy in most fields. Across the 52 canonical fields, mean coupling
is $+0.419$ \ci{+0.355}{+0.482}, and 39 of the 52 fields have coupling more than 1.96 bootstrap standard
errors above zero, none below (Fig.~\ref{fig:map}). Coupling varies across fields by more than institution
sampling alone would produce. In a three-level meta-analytic model of Fisher-$z$ coupling on 57 fields
\citep{konstantopoulos2011fixed, cheung2014modeling}, Cochran's $Q=183.2$ on 56 degrees of freedom and
$I^2=0.69$ (0.67 with the Bonett--Wright variance for Spearman correlations). These statistics count
institution sampling only. Noise in the program medians that institution resampling does not capture,
which plausibly differs by field, is counted as heterogeneity, so part of the spread may be field-varying
measurement quality rather than field-varying coupling (\SI{3}).

Part of the variation is a measurement confound. Coupling rises with the number of institutions that
offer a field (Spearman between $z_f$ and $n_f$ $=0.43$, $p<0.001$, 57 fields), which may reflect range
restriction and noisier estimates in small fields as much as anything about the fields themselves. The
grouping of fields into disciplines is therefore a descriptive map, not a statement about variance
shares. Shrunken CIP-2 family means are lowest for health (0.27), natural resources (0.28), the arts
(0.28) and the biological sciences (0.32), and highest for computer and information sciences (0.62; one
field, computer science), mathematics and statistics (0.60), the social sciences (0.56) and parks,
recreation and kinesiology (0.55). A CIP-2 variance component is detected without adjustment
(likelihood ratio 7.68, size-preserving label permutation $p=0.005$) but is borderline once field size
is a covariate (likelihood ratio 3.41; $p$ from 0.028 to 0.046 across four null distributions), and its
profile interval then reaches zero (\SI{3}).

The field ordering is similar on a second earnings source. Coupling computed from PSEO five-year earnings
correlates with Scorecard coupling across fields at $+0.68$ (two-way interval \ci{+0.54}{+0.82}; 57
fields, V4.13.0) and at $+0.74$ \ci{+0.61}{+0.87} on the V4.14.1 release (127 matched institutions). This
is robustness to the earnings measure rather than an independent replication: both couplings use the same
$F$, 61\% of the PSEO institution-by-field rows used on V4.14.1 also enter the Scorecard couplings, and
institution noise shared by the two sources could raise the correlation, which we cannot separate. Pooled
over all 1,535 matched institution-by-field cells, the two sources' earnings levels correlate at $+0.94$
(Spearman), but that figure includes pay differences between fields; within fields the mean Spearman is
$+0.68$ (median $+0.73$, range $+0.19$ to $+0.95$; 48 fields with at least 8 matched cells; SI check,
\SI{3}).

Median pay is not an idiosyncratic stand-in for where graduates go. In PSEO Flows, within CIP-2
families, the rank correlation of family prestige $P$ with the share of employed graduates in market
sectors is $+0.40$ \ci{+0.32}{+0.47}, about the same as wage coupling on the same cells, $+0.41$
\ci{+0.31}{+0.49} (1,312 institution-by-family cells at 108 institutions, 21 families; difference $+0.01$
\ci{-0.08}{+0.10}). On the 1,178 cells at 93 institutions with Scorecard covariates, where the raw values
are $+0.40$ and $+0.42$, both fall to about $+0.1$ net of selectivity (placement $+0.13$
\ci{+0.04}{+0.22}, wages $+0.10$ \ci{-0.00}{+0.21}). Across fields, graduates' earnings also track non-wage
measures of placement such as occupational prestige, although that check validates earnings as a placement
measure across fields, not within them (\SI{10}).

\begin{figure}[p]
\centering
\paperfig{paper_fig1}
\caption{Coupling between department prestige and bachelor's earnings, by US field. Rows are the 52 fields
with at least 15 matched institutions, sorted by raw coupling $\rho_f$, the Spearman correlation across
institutions between published field prestige $F$ \citep{wapman2022quantifying} and College Scorecard
median earnings four years after completion. Circles: raw $\rho_f$, coloured by CIP-2 family and sized by
the number of institutions. Open squares: partial coupling net of the broad selectivity set (institution
average SAT, admission rate, Pell share, control and state earnings level; 46 fields). Diamonds: the same
net of academia-wide prestige $G$ as well (44 fields). Field intervals are $\pm1.96$ joint
institution-bootstrap standard errors, cut at $\pm1$. The top row gives the means across fields with
joint institution-bootstrap 95\% intervals (conditional on the observed fields); vertical lines mark the
three means (dashed: raw; dotted: net of selectivity; dash-dot: net of selectivity and $G$). Bold labels
mark the 20 fields that pass the reliability flag.}\label{fig:map}
\end{figure}

\subsection{Most of the coupling is shared with institution-level status}\label{sec:status}

Controlling for institution-level status removes most of the coupling. What remains is small, borderline
on average and concentrated in a few fields. The estimand is the information in $F$ about pay net of
institution-level status proxies; it is not a return. In the framework of \citet{dale2002estimating},
institutional average SAT is a candidate quality treatment as well as a selection proxy, so partialling it
out removes selection together with any college effect correlated with it. The drop therefore measures what
cannot be separated from institutional status with public data, not what selectivity causes.

\paragraph{The control ladder.} On a common sample of 44 fields with SAT data, mean coupling falls from
$+0.43$ raw to $+0.31$ with Pell share, control and the state earnings level, to $+0.24$ with the
admission rate and to $+0.14$ with average SAT (Fig.~\ref{fig:status}a); adding institution-wide earnings
takes it to $+0.12$ (Table~\ref{tab:ladder}, panel A), although that control partly contains the outcome,
because institution-wide earnings include the field's own graduates. On the 46 fields estimable in the
broad specification (SAT, admission rate, Pell share, control and state earnings level) the mean partial
coupling is $+0.141$ \ci{+0.075}{+0.208}, against $+0.43$ raw on the same institutions. Selectivity
itself couples with field earnings more strongly than department prestige does: on the 47 fields of the
SAT sample, the mean of $\rho(\mathrm{SAT},Y)$ is $+0.57$ and that of $\rho(F,Y)$ is $+0.43$. Adding
academia-wide prestige $G$ to the broad set leaves $+0.059$ \ci{+0.003}{+0.114} (44 fields;
Fig.~\ref{fig:status}a, last step): the part of the department ordering that neither selectivity nor the
academia-wide ordering tracks. Whether this remainder excludes zero depends on the specification. Across ten
specifications that net out selectivity and $G$, the interval excludes zero in four; point estimates range
from $+0.041$ to $+0.140$, and no upper bound exceeds $+0.25$. Averaged over the four-, one- and five-year
Scorecard horizons on the 2,422 programs released at all three, it is $+0.052$ \ci{+0.003}{+0.100} (39
fields).

\paragraph{Within institutions.} Institution fixed effects absorb each institution's average earnings
level across its fields, and with it the institution-wide parts of selectivity, $G$ and location. They do
not absorb field-specific local labour demand or differences in who enrols in a given major, so the
remaining slope compares departments of different standing inside the same institution without separating
department standing from program-level selection. The slope of log four-year earnings on department
prestige falls from $+0.084$ (SE 0.006) to $+0.013$ (SE 0.004) per within-field SD when institution fixed
effects are added (3,416 programs, 198 institutions). Fixed effects do not remove field-specific slopes on
institution-level selectivity, so we add field-specific slopes on SAT, admission rate, Pell share and $G$
(Eq.~\ref{eq:fe}). The department slope is then $+0.0084$ \ci{+0.0000}{+0.0164} in log earnings (2,889
programs, 168 institutions, 47 fields; institution cluster bootstrap), against $+0.0851$ with field fixed
effects only on the same programs. About a tenth of the cross-institution slope is left: one SD of
department prestige goes with 0.84\% higher four-year pay. Its interval reaches zero only at the edge
(2.5\% of 999 bootstrap draws are at or below zero), and it survives field-specific cubic slopes in $G$
($+0.0093$ \ci{+0.0002}{+0.0176}) and a field-specific slope on the state earnings level ($+0.0097$
\ci{+0.0009}{+0.0171}), the only location control available.

The four-year outcome gives the largest remainder. With one-year earnings the main-specification slope is
$-0.0024$ \ci{-0.0150}{+0.0101} (2,681 programs) and with five-year earnings $+0.0023$ \ci{-0.0069}{+0.0119}
(2,549 programs), against $+0.0123$ and $+0.0111$ with four-year earnings on the same programs; 3 of 14
paired horizon differences exclude zero, all with a smaller remainder at one or five years. The horizons
describe different graduating cohorts, not career time. Averaged over the three horizons on the 2,422
programs released at all three (163 institutions, 42 fields), the slope is $+0.0048$ \ci{-0.0024}{+0.0120}
with each program weighted equally and $+0.0108$ \ci{+0.0035}{+0.0181} with programs weighted by the number
of graduates behind each median; the weighting changes it by $+0.0060$ \ci{+0.0005}{+0.0112}. The unweighted
estimate describes the average program; the weighted one describes the average graduate and gives noisy
small-program medians less weight. Together they put the remainder at 0.5--1.1\% of pay per SD
(Table~\ref{tab:ladder}, panel B; Fig.~\ref{fig:status}b); corrected for measurement error in department
prestige it could be larger (below).

\paragraph{Department prestige against its competitors.} In per-field rank regressions on 47 fields, the
mean Shapley shares of $R^2$ are 0.076 for $F$, 0.173 for selectivity (SAT and admission rate), 0.086
for $G$ and 0.201 for institution-wide earnings (Fig.~\ref{fig:status}c); institution-wide earnings contain
the field's own graduates and lead partly by construction, and without them the first three shares are
0.092, 0.257 and 0.107. $F$ is the largest block in 2 fields and institution-wide earnings in 28; the
difference between the shares of $F$ and $G$ is $-0.010$ \ci{-0.026}{+0.005}. Before any control, the
academia-wide ordering couples with field earnings slightly more strongly than the field's own ordering: the
mean of $c_F-c_G$, where $c_F$ and $c_G$ are the couplings of $F$ and $G$ with the same field earnings, is
$-0.049$ \ci{-0.085}{-0.011} (57 fields). $F$ is measured with more error than $G$. Correcting both for
their reliability \citep{spearman1904proof} shrinks the difference to $-0.005$ \ci{-0.043}{+0.044} at the
public-edge lower bound and to $-0.042$ \ci{-0.077}{-0.001} at the extrapolated reliability
(Fig.~\ref{fig:status}d). Net of selectivity the two are tied ($+0.00$ \ci{-0.05}{+0.05}, 46 fields). Field
prestige thus predicts pay no better than academia-wide prestige under every correction; whether it predicts
pay worse depends on its reliability, which public data cannot pin down.

\paragraph{Where the remainder sits.} Within institutions, per-field slopes differ when field-specific
slopes on selectivity and $G$ are linear (wild-cluster heterogeneity test $p\le 0.001$ over 47 fields;
$p=0.006$ with leverage-corrected residuals), and their ordering reproduces across random halves of
institutions (split-half reliability $+0.25$, permutation $p=0.02$; $p=0.06$ against a simulated null).
The heterogeneity is fragile: with field-specific cubic slopes in $G$, the test on the 23
best-estimated fields gives $p=0.09$--$0.24$ and the split-half reliability falls to $+0.05$. The one
field whose remainder is individually clear is computer science, one of three fields named in advance as
high-coupling: $+0.050$ \ci{+0.012}{+0.096} per SD pooled over horizons (128 programs), with a 98.33\%
interval, which corrects for the three fields named in advance, of \ci{+0.007}{+0.114}; with four-year
earnings it stays positive under both nonlinear variants ($+0.047$ \ci{+0.010}{+0.084} and $+0.051$
\ci{+0.012}{+0.087}). Several engineering and business fields are positive, but none passes a
multiple-testing correction across the 42 fields with percentile-bootstrap $p$-values. Two groupings were
read off the per-field results, and both are post hoc. The nine computer-science and engineering fields
average $+0.0284$ \ci{+0.0129}{+0.0432} against $-0.0006$ for the other 33; weighted by precision the group
mean is $+0.0134$ \ci{-0.0007}{+0.0271}, and random sets of nine field labels reach the equal-weight
contrast in 4.2\% of draws (11.9\% with precision weights). The five business fields average $+0.0198$
\ci{+0.0027}{+0.0354}, about as much; they exceed the remaining 28 fields by $+0.0242$ \ci{-0.0028}{+0.0462}
with equal weights and by $+0.0273$ \ci{+0.0094}{+0.0437} with precision weights, and random sets of five
labels reach the business contrast in 21.1\% of draws (7.9\% with precision weights). Refitted without the
nine computer-science and engineering fields, the pooled graduate-weighted slope is $+0.0079$
\ci{-0.0017}{+0.0161}. A coefficient-stability bound puts the selection on unobserved variables needed to
erase the within-institution slope at 0.20 to 1.34 times the selection on the observed controls, below the
benchmark of 1 in three of four settings (\SI{4}).

\paragraph{The cross-field ordering.} The ordering of fields survives adjustment in part: broad partial
couplings still rank fields like raw couplings beyond shared estimation noise (Spearman $+0.63$ against a
null mean of $+0.29$, $p=0.006$; 46 fields). But the ordering is nearly collinear with how steeply each
field's earnings follow institutional selectivity. Across the 23 fields with at least 60 SAT-reporting
institutions, raw coupling and the selectivity gradient $\rho(\mathrm{SAT},Y)$ have a true-score
correlation of $+0.94$. With institution-level proxies standing in for the quality of a program's own
students, a field-specific department reading and a field-specific selectivity-gradient reading predict
the same ordering, and public data cannot separate them (\SI{4}).

\paragraph{The upper tail.} Medians could miss status that matters in the upper tail of outcomes, where
elite-college effects concentrate \citep{chetty2023diversifying}. College-level outcomes from Opportunity
Insights show slightly tighter coupling toward the top without controls, and no extra tail loading once
selectivity and parental income are held fixed. Across 242 PhD-granting universities, $G$ tracks the share
of former students reaching the top 1\% of their cohort's earnings ($\rho=+0.66$ \ci{+0.58}{+0.74}) about
as closely as their median earnings ($+0.62$ \ci{+0.54}{+0.69}); raw, the top-5\% share exceeds the median
by $+0.04$ \ci{+0.01}{+0.08} and the top-1\% share by $+0.05$ \ci{-0.01}{+0.11}. Net of selectivity and
parental income, $G$'s partial correlation is $+0.12$ \ci{-0.00}{+0.26} with the top-1\% share and $+0.15$
\ci{+0.01}{+0.29} with the median (difference $-0.02$ \ci{-0.13}{+0.10}). A tail premium confined to the
roughly 80 most selective colleges cannot be resolved at this sample size (Barron's tiers 1--2, $n=80$: top
1\% minus median $+0.22$ \ci{+0.07}{+0.38} raw, $+0.18$ \ci{-0.07}{+0.40} controlled). These outcomes pool
all of a college's fields, and $G$ tracks the institution's average field prestige at $+0.91$, so they
cannot compare $F$ with $G$ (\SI{8}).

\paragraph{Measurement error and the two readings.} Two things bias the remainder toward zero. The residual
prestige measure is noisy, which attenuates it, and partialling $G$ also removes some genuinely
field-specific variation, because $G$ pools the same hires, the field's own included. Disattenuating with
the reliability bracket gives a mean partial coupling of $+0.159$ \ci{+0.052}{+0.266} at the lower-bound
reliability (34 of 44 fields identifiable) and $+0.064$ \ci{-0.002}{+0.129} at the extrapolated one, against
$+0.419$ raw; the within-institution slope becomes $+0.0297$ or $+0.0098$, against $+0.0851$ without
institution fixed effects. Two readings remain. Under the first, the academy's field-specific hierarchy
carries labour-market information of its own. Under the second, within-field coupling is the field's
gradient in institutional status plus whatever student quality institution-level measures miss. The bulk of
the coupling goes with institution-level status, as the second reading predicts. For the remainder the two
predict the same pattern: under the second, selection of stronger students into separately admitted or
capped majors within an institution, such as separately admitted engineering or business programs or
capped computer-science majors (restricted majors can admit students on their grades;
\citealp{bleemer2022will}), would also concentrate a premium in computer science, engineering and
business. Public data cannot tell the two apart.

\begin{figure}[!tbp]
\centering
\paperfig{paper_fig2}
\caption{What remains of US coupling after institution-level status (College Scorecard, bachelor's median
earnings four years after completion). (a) Within-field coupling of department prestige $F$ as controls are
added from left to right; grey lines are fields, the black line their mean (44 fields estimable at every
step), the bar the joint institution-bootstrap 95\% interval of the final mean. The ladder in
Table~\ref{tab:ladder} ends instead with institution-wide earnings ($+0.12$). (b) Within-institution slope,
as the change in log earnings per within-field SD of $F$, from field fixed effects only to the main
specification (institution fixed effects and field-specific slopes on SAT, admission rate, Pell share and
$G$); 95\% institution-cluster bootstrap intervals, except the grey no-fixed-effects point ($\pm1.96$
institution-clustered SE) and the black field-fixed-effects point (point estimate only). (c) Mean Shapley
shares of $R^2$ in per-field rank regressions of earnings on department prestige, selectivity (SAT and
admission rate), $G$ and institution-wide earnings (solid bars) and without the institution-earnings block
(hatched); 47 fields. (d) Mean difference $c_F-c_G$ between the couplings of field earnings with $F$ and
with $G$: uncorrected; with both couplings corrected for measurement error at the lower-bound or the
extrapolated reliability of $F$ (that of $G$ is about 1 in both); net of the broad selectivity set; and $F$
net of selectivity and $G$ minus $G$ net of selectivity and $F$. Intervals in (c) and (d): joint
institution bootstrap, 95\%.}\label{fig:status}
\end{figure}

\begin{table}[!tbp]
\centering
\caption{Coupling under institution-level controls. Panels A and B: US, College Scorecard median
earnings four years after completion unless stated, published field prestige $F$. Panel C: UK, LEO median
earnings five years after graduation (cohorts 2013/14--2016/17), academia-wide UK hiring prestige $G$.
$k$ = fields (subjects). Intervals (95\%): joint institution bootstrap, conditional on the observed fields
(A); institution cluster bootstrap (B); two-stage bootstrap over subjects and providers (C). Ladder:
point estimates on a common sample.}\label{tab:ladder}
\footnotesize
\renewcommand{\arraystretch}{0.95}
\begin{tabular}{@{}>{\raggedright\arraybackslash}p{8.1cm}>{\raggedright\arraybackslash}p{2.5cm}>{\raggedright\arraybackslash}p{4.5cm}@{}}
\toprule
Specification & $k$ / sample & Estimate [95\% interval] \\
\midrule
\multicolumn{3}{@{}l}{\emph{A. US: mean within-field coupling of $F$ with earnings}} \\
Raw, canonical sample & 52 & $+0.419$ \ci{+0.355}{+0.482} \\
Ladder: raw $\to$ + Pell share, control, state level $\to$ + admission rate $\to$ + SAT $\to$ + institution-wide earnings & 44 & $+0.43\to+0.31\to+0.24\to+0.14\to+0.12$ \\
Broad set (SAT, admission rate, Pell share, control, state level) & 46 & $+0.141$ \ci{+0.075}{+0.208} \\
Broad set + $G$ (field-specific remainder) & 44 & $+0.059$ \ci{+0.003}{+0.114} \\
Broad set + $G$ + institution-wide earnings & 44 & $+0.047$ \ci{-0.011}{+0.106} \\
Broad set + $G$, mean of 4-, 1- and 5-year earnings & 39 & $+0.052$ \ci{+0.003}{+0.100} \\
\midrule
\multicolumn{3}{@{}l}{\emph{B. US: within-institution slope, change in log earnings per within-field SD of $F$}} \\
Field fixed effects only & 3,416 programs & $+0.084$ (SE 0.006) \\
+ institution fixed effects & 3,416 programs & $+0.0129$ \ci{+0.0048}{+0.0212} \\
+ institution fixed effects, SAT sample & 2,889 programs & $+0.0139$ \ci{+0.0053}{+0.0227} \\
\quad + field-specific slopes on SAT and $G$ & 2,889 programs & $+0.0118$ \ci{+0.0037}{+0.0205} \\
\quad + field-specific slopes on SAT, admission rate, Pell share, $G$ (main) & 2,889 programs & $+0.0084$ \ci{+0.0000}{+0.0164} \\
\quad same programs, field fixed effects only & 2,889 programs & $+0.0851$ (point) \\
\quad main + field-specific cubic in $G$ & 2,889 programs & $+0.0093$ \ci{+0.0002}{+0.0176} \\
\quad main, 1-year earnings (4-year on the same programs $+0.0123$) & 2,681 programs & $-0.0024$ \ci{-0.0150}{+0.0101} \\
\quad main, 5-year earnings (4-year on the same programs $+0.0111$) & 2,549 programs & $+0.0023$ \ci{-0.0069}{+0.0119} \\
\quad main, mean of 4-, 1-, 5-year earnings: unweighted & 2,422 programs & $+0.0048$ \ci{-0.0024}{+0.0120} \\
\quad main, mean of 4-, 1-, 5-year earnings: weighted by cohort size & 2,422 programs & $+0.0108$ \ci{+0.0035}{+0.0181} \\
\midrule
\multicolumn{3}{@{}l}{\emph{C. UK: coupling of $G$ with earnings and prior attainment}} \\
Raw & 33 & $+0.486$ \ci{+0.398}{+0.559} \\
Raw, composition sample & 32 & $+0.497$ \ci{+0.406}{+0.571} \\
\quad standardised for band mix, within-provider gradients (85\% \ci{80}{90} survives) & 32 & $+0.425$ \ci{+0.335}{+0.497} \\
\quad standardised for band mix, national band medians (floor; 62\% \ci{56}{69}) & 32 & $+0.310$ \ci{+0.234}{+0.374} \\
\quad partial on institution intake selectivity (26\% \ci{17}{38}) & 32 & $+0.131$ \ci{+0.078}{+0.188} \\
Within prior-attainment band & 28 & $+0.227$ \ci{+0.160}{+0.279} \\
\quad all graduates, same providers (66\% \ci{51}{79} survives; noise-corrected 87\% \ci{74}{100}) & 28 & $+0.342$ \ci{+0.255}{+0.398} \\
\quad within band: $G$ net of intake selectivity & 28 & $+0.033$ \ci{-0.027}{+0.098} \\
\quad within band: intake selectivity net of $G$ & 28 & $+0.252$ \ci{+0.172}{+0.310} \\
\bottomrule
\end{tabular}
\end{table}

\subsection{UK: most coupling survives a direct prior-attainment control}\label{sec:uk}

The UK data add what the US data lack: earnings by graduates' own prior attainment. They have no department
axis, so they bear on institution-level status against student attainment, not on the department question.
Most UK coupling survives a control for attainment, and what survives tracks intake selectivity, not
hiring-network prestige.

UK coupling of $G$ with LEO earnings five years after graduation averages $+0.486$ \ci{+0.398}{+0.559}
across 33 subjects, similar to US coupling of academia-wide prestige with Scorecard earnings ($+0.450$, 57
fields; \SI{4}). We control for prior attainment in three ways (Table~\ref{tab:ladder}, panel C;
Fig.~\ref{fig:uk}). First, among same-band graduates across providers, coupling is $+0.227$
\ci{+0.160}{+0.279} against $+0.342$ \ci{+0.255}{+0.398} for all graduates of the same providers (28
subjects), so 66\% \ci{51}{79} survives. Band medians rest on fewer graduates (adjacent-cohort test--retest
reliability 0.48 against 0.74 for all-graduate medians), which pulls within-band coupling toward zero;
correcting for that noise gives 87\% \ci{74}{100} (27 subjects). Second, standardising each provider's
earnings for its mix of attainment bands, with the band gradient measured within providers, leaves $+0.425$
\ci{+0.335}{+0.497} of $+0.497$ \ci{+0.406}{+0.571} on the same cells: 85\% \ci{80}{90} survives (32
subjects), and 90\% \ci{84}{94} with a freely estimated gradient on the providers where it is identified.
Third, standardising with national subject-by-band medians leaves 62\% \ci{56}{69}. This over-corrects.
High-band graduates concentrate at high-paying institutions, so national band gaps mix attainment with
institution: the log earnings gap between the 360-point band and the 240--299-point band is $+0.209$
nationally against $+0.058$ within providers. We read 62\% as a floor. The bands are coarse and sorting
within a band is not controlled, which biases every share upward. The shares are ratios of rank
correlations: they say how much of the prestige ordering of pay survives the control, not how much of the
earnings gap between high- and low-$G$ providers survives, which these rank statistics do not measure.

What survives is intake selectivity. Among same-band graduates, the institution's intake selectivity
couples with earnings at $+0.326$ \ci{+0.243}{+0.387}; $G$ adds $+0.033$ \ci{-0.027}{+0.098} beyond it,
while selectivity keeps $+0.252$ \ci{+0.172}{+0.310} beyond $G$ (28 subjects). On all-graduate medians,
partialling out selectivity leaves $+0.131$ \ci{+0.078}{+0.188}, 26\% \ci{17}{38} of the raw coupling
(32 subjects), a fall comparable to the US one from $+0.43$ to $+0.14$ under different controls. This is
the pattern \citet{britton2022how} report with a composite league-table rating, found again with a
hiring-network axis. Within bands, the selectivity measure can also pick up finer sorting of students,
so the result does not separate institutional status from selection within bands. The UK $G$ may also
measure standing less validly than the US ranks (Section~\ref{sec:earnings}); validity error in $G$ would
shift weight in these partial correlations from $G$ to the correlated selectivity measure, so the $+0.033$
may understate what $G$ carries.

Subjects differ; the differences survive the attainment control but shrink and reorder once intake
selectivity is partialled out. The between-subject SD of true coupling is 0.16 for all graduates
($I^2=0.84$) and 0.17 after within-provider standardisation ($I^2=0.83$; 32 subjects), and the subject
ordering barely moves (Spearman $+0.94$ \ci{+0.89}{+0.98}); without the three pay-scale subjects of
Section~\ref{sec:boundary} the SD is 0.12 for all graduates and 0.14 after standardisation. But subjects
whose pay couples strongly with $G$ are the subjects whose pay couples strongly with intake selectivity
(Spearman across subjects $+0.92$ \ci{+0.71}{+0.95}). Net of selectivity, less variation remains and it is
ordered differently (SD 0.09; Spearman with the raw ordering $+0.31$ \ci{+0.06}{+0.66}). As in the US, the
subject ordering cannot be separated from subject-specific selectivity gradients. Finally, on the 16
subjects where a field-tagged UK prestige axis is usable, $G$ couples with pay more strongly than the field
axis in all 16 (mean $+0.639$ against $+0.417$ on the same providers). The field-tagged axis rests on far
fewer hires, so this shows that field prestige is no better than $G$, not that it is worse in truth.

\begin{figure}[!tbp]
\centering
\paperfig{paper_fig3}
\caption{UK: coupling and graduates' prior attainment. LEO median earnings five years after graduation,
cohorts 2013/14--2016/17; academia-wide UK hiring prestige $G$ from ORCID-derived PhD-to-faculty moves. (a)
Coupling by CAH2 subject for all graduates (filled) and after standardising each provider's earnings for
its prior-attainment mix with within-provider band gradients (open), on the composition sample of 32
subjects (providers with attainment-mix data); the top row gives the means with two-stage 95\% intervals.
Because the sample differs, subject values for all graduates differ from Fig.~\ref{fig:boundary}b and the
``Raw'' column of the per-subject table in \SI{5}, which use all 134 providers (for example, education and
teaching $+0.097$ here and $+0.220$ there). Orange diamonds mark the three subjects largely on national pay scales (filled
label symbol) and allied health (open label symbol). (b) Share of coupling that survives each
prior-attainment control, with 95\% two-stage bootstrap intervals: same-band against all graduates of the
same providers; the same with the band-noise correction; within-provider standardisation; a free
within-provider band gradient (subset of providers); national band medians, which over-correct (hatched).
(c) Coupling of $G$ and of intake selectivity (the institution's share of graduates with at least 360 UCAS
points) with earnings, and each net of the other ($A\,|\,B$: partial correlation of $A$ net of $B$), for
same-band graduates (28 subjects); for all graduates of the same providers, $\rho(G)$, $\rho$(selectivity)
and $G$ net of selectivity.}\label{fig:uk}
\end{figure}

\subsection{Coupling rises with years since graduation}\label{sec:career}

Within fixed graduation cohorts, coupling rises with years since graduation in both countries. In the US the
rise appears in the loading on academia-wide prestige $G$; the loading on field prestige $F$ does not grow.

\paragraph{US.} On the balanced PSEO panel, mean coupling rises from $+0.173$ one year after graduation to
$+0.310$ at five years and $+0.451$ at ten (39 fields, V4.13.0; Fig.~\ref{fig:career}a,
Table~\ref{tab:career}). The mean per-field slope is $+0.031$ per year (two-way interval
\ci{+0.021}{+0.041}; crossed bootstrap \ci{+0.018}{+0.041}); coupling of $G$ rises from $+0.151$ to
$+0.474$, by $+0.036$ per year (two-stage interval \ci{+0.027}{+0.043}). The rise is positive in each of the
four cohorts and in both segments (y1 to y5 $+0.034$ \ci{+0.015}{+0.053}; y5 to y10 $+0.028$
\ci{+0.020}{+0.037}), so it does not rest on the first-year point. Several fields have negative first-year
coupling, most of all kinesiology and music ($-0.35$), environmental sciences ($-0.19$) and anthropology,
social work and geography ($-0.15$). This is consistent with, though not a measurement of, graduates of
higher-prestige institutions more often being in graduate school at y1 (\SI{10}).

A within-cohort change mixes career time with calendar time. Comparing cohort $c$ at y10 with cohort $c+9$
at y1, which fall in about the same calendar years, gives $+0.014$ per year (two-way \ci{+0.007}{+0.022}),
and the drift across cohorts at y1 is $+0.017$ \ci{+0.010}{+0.025}. Under a linear age--period--cohort
reading, the within-cohort change, the calendar-matched contrast and the drift identify $a+b$, $a-g$ and
$b+g$, where $a$, $b$ and $g$ are career, calendar and cohort effects (Section~\ref{sec:methods}). If
calendar and cohort effects are both non-negative, the career-time part of the rise lies between $+0.014$
and $+0.032$ per year \citep{fosse2019bounding}. Without that assumption it is not identified, and the
assumption carries weight: a zero career-time effect would need a calendar trend of about $+0.032$ per year
together with a cohort trend of about $-0.014$ per year, which we cannot rule out over the years the panel
spans, 2002 to 2022, which include the Great Recession and the pandemic (the y10 earnings of the 2010--12
cohort fall in 2020--22). The calendar-matched contrast is similar across the four cohort pairs ($+0.009$ to
$+0.016$ per year), but the pairs share calendar years. The August 2026 release changes no statement
($+0.030$ per year, two-way \ci{+0.020}{+0.041}, 40 fields; bound $[+0.015, +0.032]$), but this is a weak
test, because the new release barely changes the fixed-cohort panel.

Two checks bear on selection into the earnings data. Higher-prestige institutions have a smaller share
of graduates with PSEO earnings at y1 (mean within-cell Spearman $-0.165$); partialling coverage out of
coupling removes 0.19 \ci{0.09}{0.29} of the slope, while restricting to high-coverage institutions
removes none \citep[cf.][]{foote2026attrition}. Coverage counts graduates with covered earnings; neither
check addresses the mix of later degrees among those covered. We detect no difference between high- and
low-coupling fields. Fields with high baseline coupling (Scorecard coupling at least 0.45 and passing the
reliability flag; 14 fields) rise by $+0.008$ per year more than the others (two-way \ci{-0.010}{+0.026}),
and across fields the slope is unrelated to baseline coupling (Spearman $+0.02$, $p=0.89$). The data exclude
a difference as large as the all-field slope but not one half that size (one-sided bound $+0.022$ per year;
power 42\% against half the slope; \SI{9}).

\paragraph{Graduate training.} Graduates of higher-status institutions may more often go on to doctoral,
medical, law or business programmes. That would lower their covered earnings early and add
graduate-degree returns later, raising coupling within a cohort without any change in how the bachelor's
labour market treats status; because the propensity would be institution-wide, it would also load on $G$
and hold $b_G$ near zero at y1 (below). An exploratory field-level check points that way for the late
segment. Across the 27 fields with a graduate-degree share (New York Fed college labour-market data, legacy
\texttt{scripts/35b}), the y5-to-y10 change in coupling correlates with that share at Spearman $+0.53$
($p=0.005$; field bootstrap \ci{+0.08}{+0.84}) and the whole slope at $+0.30$ \ci{-0.17}{+0.73}. The 11
fields where fewer than 40\% of graduates take a graduate degree rise by $+0.021$ per year ($+0.015$ in the
late segment) against $+0.031$ ($+0.037$) in the 8 fields where half or more do: the rise is present in
mostly terminal-bachelor's fields, but smaller. The growth of $b_G$ does not track the share (Spearman
$-0.05$ \ci{-0.47}{+0.40}) and is about the same in the two groups ($+0.037$ and $+0.038$ per year). These
are SI-check computations on the field table of \texttt{scripts/52}: correlations across fields that ignore
shared institutions and have no two-way interval. PSEO cannot separate the returns to later degrees from
career dynamics in the bachelor's labour market.

\paragraph{Academia-wide against field prestige.} In a standardised rank regression of earnings on $F$ and
$G$ within each field-by-cohort cell, the field-prestige coefficient $b_F$ stays flat ($+0.189$, $+0.176$
and $+0.177$ at y1, y5 and y10), while the academia-wide coefficient $b_G$ grows from $-0.013$ at y1 to
$+0.328$ at y10 (Fig.~\ref{fig:career}c). Per year, $b_G$ rises by $d_G=+0.038$ (two-way
\ci{+0.021}{+0.055}) and $b_F$ by $d_F=-0.001$ \ci{-0.016}{+0.013}; the difference $d_G-d_F$ is $+0.039$
\ci{+0.009}{+0.069} ($z=2.57$) on V4.13.0 and $+0.035$ \ci{+0.007}{+0.063} ($z=2.45$) on V4.14.1, and above
zero under the two-way variance in all five data variants we checked ($z$ from 2.45 to 2.65). The evidence is
moderate, not decisive. In a calibration on permuted data the two-way variance of $d_G-d_F$ is 1.04 times the
true sampling variance on average but 0.54 to 1.71 times on single datasets; the conservative crossed
bootstrap reaches zero (\ci{-0.003}{+0.070} on V4.14.1); the five variants are nested versions of almost the
same panel; and the two-way standard replaced the crossed bootstrap during the analysis, after the calibration
showed the latter's variance to be 3.0 to 3.7 times the true sampling variance (an earlier revision of
\texttt{scripts/59} had rated the difference clear only with fields held fixed; \SI{6}). The decomposition is
also weakly identified: within cells $F$ and $G$ correlate 0.83 and $G$ is measured more reliably, so least
squares shifts shared signal toward $G$ by an amount public data do not identify, and we have not re-estimated
the loadings with $F$ corrected for measurement error. At y1 the association runs through $F$ ($b_F=+0.189$,
two-way \ci{+0.062}{+0.317}; $b_G=-0.013$ \ci{-0.160}{+0.134}). If $F$ were only a noisier copy of $G$, $G$
would take the larger coefficient at every horizon, so the first-year pattern is consistent with more than one
component; but y1 is also the horizon most affected by coverage and graduate study. The regression does not
say what the growing $G$-tracked component is: institutional reputation, selectivity ($G$ correlates $+0.68$
with average SAT), resources, location or the timing of graduate training.

\paragraph{Quantiles.} Coupling of the 25th percentile of earnings is below median coupling by $-0.067$
(two-way \ci{-0.092}{-0.042}; V4.13.0, mean over horizons). Coupling of the 75th percentile minus median
coupling is $+0.027$ (two-way \ci{+0.001}{+0.053} on V4.13.0 and \ci{+0.000}{+0.054} on V4.14.1); the
crossed intervals include zero on both releases and the institution-clustered permutation at y10 gives
$p=0.34$--$0.37$, so we do not claim a difference. The robust pattern is at the bottom: the lower quartile
is less ordered by prestige than the median. PSEO quartiles are within-program quantiles, not the elite
tail.

\paragraph{UK.} On fixed LEO cohorts (2013/14 to 2016/17, providers observed at YAG 1, 3 and 5; 33
subjects), coupling rises from $+0.379$ at YAG1 to $+0.484$ at YAG5, $+0.026$ per year
\ci{+0.014}{+0.037} (permutation $p=0.001$; Fig.~\ref{fig:career}b). It is positive in all four cohorts
and in both segments (YAG1 to YAG3 $+0.029$ \ci{+0.007}{+0.047}; YAG3 to YAG5 $+0.023$
\ci{+0.009}{+0.039}), and it survives the prior-attainment controls: $+0.028$ \ci{+0.016}{+0.038} per year
after within-provider standardisation and $+0.017$ \ci{+0.002}{+0.031} within bands, the latter
attenuated by band-median noise. A coverage control lowers it to $+0.017$ \ci{+0.005}{+0.027}. The
calendar-matched contrast (cohort $c$ at YAG5 against cohort $c+4$ at YAG1, the same tax year) is
$+0.031$ \ci{+0.021}{+0.042}, the within-cohort change on the same providers $+0.025$
\ci{+0.015}{+0.036}, and the cross-cohort drift at YAG1 $-0.006$ \ci{-0.015}{+0.003}. These point
estimates imply that calendar and cohort effects sum to slightly below zero, so the sign assumption of the
US bound holds here only within sampling error, and offsetting calendar and cohort effects cannot be
excluded. The pandemic and furlough tax years 2020/21 and 2021/22 fall both in the YAG5 cells of the fixed
cohorts and in the YAG1 cells of the calendar contrast.

\paragraph{Reading.} The direction agrees with descriptive US and Colombian evidence
\citep{thomas2005postbaccalaureate, macleod2017big}. It measures a different estimand from causal estimates
of small or fading premia \citep{bordon2020employer, mountjoy2021returns} and does not contradict them: a
descriptive status gradient can rise while a selection-corrected premium fades. Coupling is unconditional,
and no public program-level measure of student ability exists, so these data cannot test employer learning
(Section~\ref{sec:discussion}). A rising association between status and pay fits a growing return to
status; it also fits graduates of higher-status institutions sorting into careers with steeper earnings
profiles \citep{binder2016career, rivera2015pedigree}, and the timing of graduate training. One descriptive
pointer: on the fixed 2010--12 PSEO cohort, the prestige ordering of market-sector placement is already
present at y1 ($\rho(P,M)=+0.28$, $+0.38$ at y10), while wage coupling rises from $+0.21$ to $+0.51$ (point
estimates, one cohort). In this cohort, sorting by status into sectors is visible before most of the growth
in the pay gradient.

\begin{figure}[!tbp]
\centering
\paperfig{paper_fig4}
\caption{Coupling and years since graduation on fixed cohorts. (a) US: mean within-field coupling of
department prestige $F$ and of academia-wide prestige $G$ with PSEO median earnings one, five and ten
years after graduation; balanced fixed cohorts of 2001--03 to 2010--12 graduates (39 fields, release
V4.13.0); bars are crossed field-by-institution bootstrap 95\% intervals. $F$'s coupling rises about as
steeply as $G$'s partly through its overlap with $G$ (within-cell Spearman 0.83); panel (c) separates the
two. (b) UK: mean coupling of academia-wide UK prestige $G$ with LEO median earnings one, three and five
years after graduation, for all graduates (33 subjects), after standardising for prior-attainment mix (32)
and within prior-attainment bands (27); balanced fixed cohorts 2013/14--2016/17; bars are two-stage
subject-and-provider bootstrap 95\% intervals. (c) US: mean standardised rank-regression coefficients of
earnings on $F$ ($b_F$) and $G$ ($b_G$) by horizon, same panel as (a); bars are two-way (field, institution)
cluster 95\% intervals; $d_F$ and $d_G$ are their slopes per year. Slopes and their intervals are in
Table~\ref{tab:career}. Within-cohort changes mix career time with calendar time
(Section~\ref{sec:career}).}\label{fig:career}
\end{figure}

\begin{table}[!tbp]
\centering
\caption{Career-time summary. US: PSEO bachelor's graduates, balanced fixed cohorts 2001--03 to
2010--12, field prestige $F$; per-field slopes averaged over fields; two-way (field, institution) cluster
intervals unless marked. UK: LEO, balanced fixed cohorts 2013/14--2016/17, academia-wide prestige $G$;
two-stage bootstrap over subjects and providers. Slopes are changes in coupling per year since
graduation. 95\% intervals.}\label{tab:career}
\footnotesize
\setlength{\tabcolsep}{4pt}
\begin{tabular}{@{}>{\raggedright\arraybackslash}p{4.3cm}lll@{}}
\toprule
Quantity & US, V4.13.0 (39 fields) & US, V4.14.1 (40 fields)$^{a}$ & UK (33 subjects) \\
\midrule
Coupling at first / middle / last horizon & $+0.173$ / $+0.310$ / $+0.451$ & $+0.159$ / $+0.294$ / $+0.433$ & $+0.379$ / $+0.438$ / $+0.484$ \\
Within-cohort slope & $+0.031$ \ci{+0.021}{+0.041} & $+0.030$ \ci{+0.020}{+0.041} & $+0.026$ \ci{+0.014}{+0.037} \\
\quad crossed bootstrap (conservative) & \ci{+0.018}{+0.041} & \ci{+0.018}{+0.043} & --- \\
Slope of $G$ coupling, two-stage interval & $+0.036$ \ci{+0.027}{+0.043} & --- & (as above) \\
Early segment (y1--y5; YAG1--3) & $+0.034$ \ci{+0.015}{+0.053} & $+0.034$ \ci{+0.016}{+0.052} & $+0.029$ \ci{+0.007}{+0.047} \\
Late segment (y5--y10; YAG3--5) & $+0.028$ \ci{+0.020}{+0.037} & $+0.028$ \ci{+0.018}{+0.037} & $+0.023$ \ci{+0.009}{+0.039} \\
Calendar-matched contrast & $+0.014$ \ci{+0.007}{+0.022} & $+0.015$ \ci{+0.007}{+0.023} & $+0.031$ \ci{+0.021}{+0.042} \\
Cross-cohort drift at entry & $+0.017$ \ci{+0.010}{+0.025} & $+0.018$ \ci{+0.010}{+0.025} & $-0.006$ \ci{-0.015}{+0.003} \\
Career-time bound if calendar, cohort effects $\ge 0$ & $[+0.014, +0.032]$ & $[+0.015, +0.032]$ & --- \\
Coverage control: share of slope removed (US); slope net of coverage (UK) & $0.192$ \ci{0.089}{0.294} & $0.195$ \ci{0.097}{0.293} & $+0.017$ \ci{+0.005}{+0.027} \\
Prior-attainment standardised slope & --- & --- & $+0.028$ \ci{+0.016}{+0.038} \\
High-coupling minus other fields & $+0.008$ \ci{-0.010}{+0.026} & $+0.006$ \ci{-0.011}{+0.024} & --- \\
Slope of $b_G$ ($d_G$) & $+0.038$ \ci{+0.021}{+0.055} & $+0.035$ \ci{+0.019}{+0.052} & --- \\
Slope of $b_F$ ($d_F$) & $-0.001$ \ci{-0.016}{+0.013} & $+0.000$ \ci{-0.013}{+0.014} & --- \\
$d_G-d_F$ (two-way $z$: 2.57; 2.45) & $+0.039$ \ci{+0.009}{+0.069} & $+0.035$ \ci{+0.007}{+0.063} & --- \\
\quad crossed bootstrap (conservative) & \ci{-0.000}{+0.078} & \ci{-0.003}{+0.070} & --- \\
$b_F$ at y1 & $+0.189$ \ci{+0.062}{+0.317} & $+0.157$ \ci{+0.050}{+0.264} & --- \\
$b_G$ at y1 / y10 & $-0.013$ / $+0.328$ & $+0.009$ / $+0.327$ & --- \\
p25 minus p50 coupling (mean over horizons) & $-0.067$ \ci{-0.092}{-0.042} & $-0.066$ \ci{-0.091}{-0.042} & --- \\
\bottomrule
\end{tabular}

\smallskip
\begin{minipage}{0.97\textwidth}\footnotesize
$^{a}$ The calendar-matched contrast, the drift and the bound use the 39 fields with triple-matched
institutions. The p75-minus-p50 difference is not distinguishable from zero and is reported in \SI{6}. The
slope of $G$ coupling is from \texttt{scripts/52}, which reports only a two-stage interval.
\end{minipage}
\end{table}

\subsection{Where coupling is weak: national pay scales and strictly licensed occupations}\label{sec:boundary}

The prestige--pay association is weak in UK subjects paid on national scales and, across US fields, falls
with the share of graduates in strictly licensed occupations. We treat both as boundary conditions of the
main finding: contrasts between discipline families that public data cannot tie to how pay is set. Special
education, whose graduates mostly teach on K-12 salary schedules, is a clear counter-case.

\paragraph{UK.} Nursing and midwifery, medicine and dentistry, and education and teaching, three subjects
whose graduates are largely paid on national scales, couple at $+0.120$ \ci{-0.033}{+0.258} against
$+0.523$ \ci{+0.439}{+0.585} for the other 30 subjects (difference $-0.404$ \ci{-0.566}{-0.238}; exact
subject-label permutation $p=0.0009$; Fig.~\ref{fig:boundary}b). After within-provider standardisation for
prior attainment the contrast is $+0.053$ against $+0.463$ ($p=0.0014$) on the 32-subject composition
sample, where the unstandardised means are $+0.057$ and $+0.543$. Allied health, much of it also paid on
national bands, is an exception at $+0.570$; adding it to the group shrinks the difference to $-0.289$
\ci{-0.497}{-0.037} ($p=0.008$). Teaching also differs over time, rising from $+0.012$ at YAG1 to $+0.211$
at YAG5. The pattern, and its reading as centrally regulated pay for medicine and education, is
\citeauthor{britton2022how}'s \citeyearpar{britton2022how}; we add nursing and a hiring-network axis.

\paragraph{US.} Net of geography, how strongly department prestige is associated with pay falls with the
share of a field's full-time workers with a bachelor's degree or more who are in strictly licensed
occupations (health practitioners, school teachers, lawyers, architects and psychologists; ACS 2023): the
Spearman correlation across fields is $-0.53$ to $-0.48$ over two institution samples and three geography
codings (Fig.~\ref{fig:boundary}a). Its interval excludes zero in all 6 versions when fields are resampled,
but in only 3 of 6 (the full institution sample) when whole CIP-2 families are resampled, and the
association rests on the ordering of discipline families. The 15 lowest strict shares among the 40 fields
of the selectivity sample all belong to engineering, computer science, statistics or business, families
whose geography-netted association is high (family means $+0.52$ to $+0.65$). Only three fields have a
strict share above 0.5: nursing (0.77), whose partial correlation net of region fixed effects is $+0.28$
\ci{+0.08}{+0.46}; communication disorders (0.61); and special education (0.61), whose graduates mostly teach on
district salary schedules and whose field couples strongly (raw $+0.55$; net of geography $+0.59$ to $+0.66$;
19 institutions), against the licensing reading. For raw coupling the rank association is weaker
($-0.29$ \ci{-0.56}{+0.04} \cic{-0.59}{+0.06}; 40 fields). It is not detected within CIP-2 families ($-0.11$ to
$+0.19$), a low-power test because only 15--21\% of the share's rank variance lies within families, nor with a
broad licensure share that adds licensed engineers, accountants, health technicians and counsellors
($-0.22$ to $-0.07$). It keeps its sign net of selectivity ($-0.42$ \ci{-0.68}{-0.10} \cic{-0.75}{-0.02}), but
licensed fields also have flatter selectivity gradients, so the pattern cannot distinguish a looser
department ordering from a flatter status gradient in general; and a field's government-or-non-profit
employment share predicts it at least as well ($-0.62$ to $-0.44$), so licensure and employer sector are not
separated.

Within fields, in a hierarchical model of program earnings with field-specific slopes on selectivity and
$G$, a field whose strict share is one SD higher (rank form) has a flatter prestige slope, $-0.088$
\ci{-0.147}{-0.049} in normal-score units against an average slope of $+0.083$ (field-label permutation
$p=0.003$, a test that ignores CIP-2 clustering). The moderation survives controls for institution-wide
earnings and institution effects but cannot be separated from field type: with an indicator for
engineering, computer science, mathematics and business fields as a second moderator it is $-0.065$
($p=0.094$), and within the 23 other fields $-0.038$ \ci{-0.127}{+0.043}. Neither the broad share ($-0.008$
\ci{-0.052}{+0.058}) nor an ex-ante licensed-field indicator ($+0.056$, $p=0.57$; 3 licensed fields) moderates
the slope in the same specification (\SI{7}). Once selectivity is netted as well as geography, nursing's
association is about zero ($-0.07$ \ci{-0.28}{+0.15}), like biology's ($-0.07$) and like that of most fields
(35 of 46 intervals include zero).

\paragraph{Employer sector in PSEO Flows.} Across institution-by-family cells, the slope of log earnings on
standardised SAT is lower by $0.016$ \ci{0.010}{0.022} for each 10 percentage points of graduates working in
education, health care and public administration (4,773 cells, 348 institutions); department prestige and $G$
show no robust flattening ($-0.004$ \ci{-0.012}{+0.004} and $-0.006$ \ci{-0.015}{+0.003}). Whether the SAT
flattening goes with this setting share or with the market share depends on the design, that is, on which
cohorts' shares are set against which cohorts' pay: shares measured at y5 load on the market share (setting
share given market share $-0.004$ \ci{-0.010}{+0.002}; market given setting $+0.009$ \ci{+0.002}{+0.016}), y1
shares pooled over cohorts load on the setting share, and each cohort's own y1 shares set against its own y5
pay load on the market share or on neither (\SI{7}). The attribution is not identified.

Pay set by the setting, such as a teacher salary schedule \citep{biasi2021labor}, an employer's pay scale
or a reimbursement rate, is a reading consistent with these contrasts and with
\citeauthor{britton2022how}'s interpretation of UK medicine and education; whether licensure itself raises
wages is contested \citep{kleiner2013analyzing, redbird2017new}. These data do not isolate either reading.
Fields with a higher licensed share also do not show less cross-institution pay dispersion (correlation
$+0.02$ \ci{-0.30}{+0.32}, 45 fields; legacy \texttt{scripts/35c}), an interval too wide to say whether pay
varies less across institutions in these fields.

\begin{figure}[!tbp]
\centering
\paperfig{paper_fig5}
\caption{Where coupling is weak. (a) US: partial correlation of department prestige $F$ with College
Scorecard four-year earnings net of BEA-region fixed effects (full institution sample, one of the six
sample-by-geography versions in Section~\ref{sec:boundary}; 95\% intervals), against the share of a
field's full-time workers with a bachelor's degree or more who are in strictly licensed occupations
(health diagnosing and treating practitioners, preschool to secondary and special-education teachers,
lawyers and judges, architects, psychologists; ACS 2023). One point per field, coloured by CIP-2 family and sized by the number
of institutions; computer science, accounting, biology, special education, nursing and communication
disorders are labelled. Top: the Spearman correlation across fields with its field-bootstrap interval and
its CIP-2 family-bootstrap interval, the within-family association (37 fields), and the association with the
broad licensure share (interval \ci{-0.50}{+0.13}). The association is a contrast between families, is not
detected within them, and is not detected with the broad share. (b) UK: coupling of academia-wide prestige
$G$ with LEO earnings five years after graduation, all graduates, full sample of 33 subjects: the three
subjects largely on national pay scales (nursing and midwifery, medicine and dentistry, education and
teaching; filled diamonds) and the other 30 (allied health, also largely on national pay bands, open
diamond); group means with 95\% two-stage bootstrap intervals. After within-provider standardisation for
prior-attainment mix, on the 32-subject composition sample of Fig.~\ref{fig:uk}a (3 and 29 subjects), the
group means are $+0.053$ and $+0.463$ ($p=0.0014$; unstandardised on that sample $+0.057$ and $+0.543$).}\label{fig:boundary}
\end{figure}

% =====================================================================================
\section{Discussion}\label{sec:discussion}

\paragraph{For the science of science.} The faculty-hiring hierarchy is a reliable measure of standing
inside the academy, the academia-wide ordering especially, and it does track the pay of the academy's
bachelor's graduates. Outside academia, however, the information it carries about their pay is mostly
institution-level status. Within institutions, one SD of department prestige goes with about 1\% higher pay
once each field's slopes on selectivity and academia-wide prestige are allowed for (0.84\%, 95\% interval
0.0--1.6\%, for four-year earnings; 0.5--1.1\% averaged over horizons; up to about 3\% after correcting for
measurement error at the lower reliability bound; Section~\ref{sec:status}); a field's own ordering
predicts pay no better than the academia-wide ordering in either country; and in the UK the academia-wide
ordering adds little beyond intake selectivity (Section~\ref{sec:uk}). The two orderings are nested, since
$G$ pools the same hiring data over all fields, so ``institution-level status'' here is the part of the
academy's hierarchy shared across an institution's departments. Graduate pay lines up mostly with that
shared part, and only a little with what distinguishes a department from its neighbours in the same
institution: about 1\% of pay per SD, individually clear only in computer science. Any study that uses
hiring-network prestige as a proxy for the training quality or labour-market standing of a program should
expect it to behave mostly as a measure of institutional status. The one field-specific exception we can
establish, computer science, is a single field, and its premium fits selection into the major as well as
department standing.

\paragraph{For the college-quality literature.} A peer-defined, hiring-network axis reproduces the UK
finding that academic quality adds little beyond selectivity \citep{hussain2009university,
britton2022how}, and extends it to about 50 US fields on College Scorecard, with the raw field ordering
reproduced on PSEO. The cross-field ordering of coupling is real, but it is nearly collinear with how
steeply each field's pay follows institutional selectivity (true-score correlation $+0.94$ in the US;
Spearman $+0.92$ across UK subjects); \citet{bloem2024understanding} show on an earlier Scorecard release
that this selectivity gradient itself varies across fields. At the resolution of public data, ``the status
gradient in pay differs by field'' and ``the department-prestige gradient in pay differs by field'' cannot
be told apart. The within-institution design narrows the gap between them to a small remainder, which is
the quantity that future data must resolve.

\paragraph{Career time.} Within fixed cohorts, coupling rises by about $+0.03$ per year in two countries,
with different earnings sources, prestige measures and horizons; in the US, in a weakly identified
decomposition, the rise appears in the loading on the academia-wide ordering, which tracks selectivity.
Employer-learning models do not predict that an unconditional association of a credential with pay falls:
in them the weight on an easily observed credential falls relative to hard-to-observe correlates of
productivity \citep{altonji2001employer}, while the unconditional association of schooling with wage levels
need not change \citep{farber1996learning}. A rising unconditional coupling fits learning if $G$ carries
productivity information that employers do not observe at entry. It also fits sorting into career tracks with
steeper earnings profiles, with which the one-cohort PSEO Flows sequence is consistent, a growing return to
status, and the timing of graduate training, for which the field-level check in Section~\ref{sec:career}
gives some support in the late segment. The data do not distinguish these. What they do establish is
descriptive: in the US the rise appears in every cohort and in both segments, it survives a coverage control
in both countries and a prior-attainment control in the UK, and in the US its career-time part lies between
$+0.014$ and $+0.032$ per year if calendar and cohort effects are non-negative. Whether it differs between
high- and low-coupling fields is not detected, by a test that can exclude only large differences.

\paragraph{What would decide the open questions.} Four kinds of data would separate department standing
from program-level selection. First, program-level entry attainment: no public US file has it, and it
would test whether the computer-science, engineering and business remainders survive students' own
attainment, as the UK band data allow for institutions. Second, changes in department standing over time:
pairing published department ratings from different assessment rounds (such as the National Research
Council assessments of 1995 and 2010) with PSEO graduation cohorts would allow a within-department design
that asks whether pay follows a department's standing when its institution's selectivity is held fixed.
Third, the admission route into majors: whether an institution admits students directly into computer
science, engineering or business, or caps entry, can be coded from public admissions pages; under selection
the premium should sit where entry to the major is selective, and under department standing it should not
depend on the route. Fourth, linked student records with entry scores, major and earnings would allow
comparisons of students with the same attainment across departments of different standing; they are the
decisive design and are not public. For career time, graduate-degree attainment by institution and field,
and selectivity measured for the entry years of each earnings cohort, would test the graduate-training
reading and let the academia-wide loading be split from selectivity.

% =====================================================================================
\section{Limitations}\label{sec:limits}

Each of the following limits a specific claim. \emph{Medians and coverage.} Scorecard earnings are medians
for Title IV--aided completers, a selected subset of graduates at elite private institutions, which may
weaken measured coupling at the top of the hierarchy; PSEO, which covers all graduates with UI-covered
earnings, lacks most top-quartile private institutions (20 of 93 in the top quartile) and cannot check this.
Neither source sees the upper tail within a program, and the Opportunity Insights check cannot resolve the
roughly 80 most selective colleges. \emph{Selectivity proxies.} Selectivity is measured at the institution
level and from the ``most recent'' Scorecard institution file: average SAT (of submitters only, in a partly
test-optional era) and admission rates of recent entering classes, while the earnings describe completers who
entered about a decade earlier (Scorecard) or in the late 1990s and 2000s (PSEO). Historical institution files
matched to each cohort's entry years would remove this mismatch; we have not used them. Whatever the controls
miss stays in every partial coupling, is correlated with prestige and weighs more in fields with steeper
selectivity gradients. \emph{Location.} No control captures field-specific local demand beyond a
field-specific slope on the state earnings level. \emph{Earnings reliability.} How reliably program medians
measure a program's pay plausibly differs by field (Section~\ref{sec:map}), and this is not corrected.
\emph{Timing.} The career-time estimates cannot separate career dynamics from calendar time without a sign
assumption, or from the timing and returns of graduate degrees. \emph{Prestige reliability and validity.} The
reliability of the published field ranks can only be bracketed (0.81--0.96), so the size of the department
remainder is uncertain: disattenuated, it ranges from $+0.064$ to $+0.159$ in partial coupling and from
$+0.0098$ to $+0.0297$ in the within-institution slope; the UK hierarchy rests on self-reported ORCID careers.
\emph{Rank statistics.} Coupling and the UK survival shares describe orderings, not the size of pay
differences; low coupling in UK pay-scale subjects may reflect little cross-provider pay variation as much as
a flat status gradient. \emph{Descriptive design.} Nothing here separates value added, selection and employer
screening; the attenuation of coupling under controls measures overlap, not mechanism. \emph{Field-level
inference.} Moderators measured once per field rest on 40--50 fields and, when families are resampled, on
about 18 CIP-2 families. \emph{Power of the nulls.} The nulls reported here are statements of what the data
cannot detect, not evidence of absence; for example, the high- against low-coupling career contrast has 42\%
power against half the all-field slope. \SI{9} gives the power and bounds of every null, including three
field-level nulls from earlier analyses of this project (a within-occupation dispersion law, occupational
concentration and disciplinary kinship). \emph{Beliefs.} Students' and employers' beliefs about prestige are
not measured. \emph{Specification.} The canonical specification (published ranks, four-year Scorecard medians,
fields with at least 15 institutions) was fixed before the manuscript was drafted, after the analyses had been
run; \SI{4} shows how the main statistics move with the alternatives, and \SI{9} lists the hypotheses examined.

% =====================================================================================
\section{Methods}\label{sec:methods}

\paragraph{Within-institution model.} For the program of field $f$ at institution $i$,
\begin{equation}
\log Y_{if}=\alpha_i+\gamma_f+\beta\,z(F)_{if}+\sum_{f'} D_{f'}\big(d_{f'}\,z\mathrm{SAT}_i+a_{f'}\,z\mathrm{ADM}_i+p_{f'}\,z\mathrm{Pell}_i+t_{f'}\,zG_i\big)+e_{if},\label{eq:fe}
\end{equation}
where $z(\cdot)$ standardises within field, $\alpha_i$ and $\gamma_f$ are institution and field fixed
effects, and $D_{f'}$ indicates field $f'$. The field-specific slopes let the earnings gradients in
selectivity, Pell share and $G$ differ by field, so $\beta$ is identified only from departments whose
standing differs from what their institution's traits imply. Institutions with one field and fields with
fewer than 15 programs are dropped iteratively. Intervals come from an institution cluster bootstrap with
999 draws, shared across specifications on the same programs so that differences are paired (\SI{4}).

\paragraph{Reliability bracket.} The published ranks cannot be re-estimated from public data, because the
public edge list omits institution pairs linked by a single hire, so we bracket their sampling reliability.
For the lower bound, each person on a public edge is assigned at random to one of two halves, SpringRank is
estimated on each half, and the Spearman correlation of the two rankings over a field's analysis
institutions is averaged over 200 seeded splits and stepped up with the Spearman--Brown formula; this treats
the published rank as if it used only the public hires. For the upper end, the same reliability is
extrapolated to the number of hires behind the published ranks with $\mathrm{rel}(m)=m/(m+m_0)$, where
$m_0$ is fitted to the public split-half. Couplings are disattenuated by the square root of the prestige
reliability; the earnings side is not corrected (\SI{2}).

\paragraph{Career-time contrasts.} On institutions observed in all three cells, the within-cohort change
$(\text{y10}-\text{y1})/9$, the calendar-matched contrast (cohort $c$ at y10 against cohort $c+9$ at y1,
divided by 9) and the cross-cohort drift at y1 (cohort $c+9$ against cohort $c$, divided by 9) equal
$a+b$, $a-g$ and $b+g$ in a linear model of coupling with career, calendar and cohort effects $a$, $b$ and
$g$. The career-time effect $a$ is bounded only under $b,g\ge0$ \citep{fosse2019bounding}; no point estimate
of it is reported (\SI{6}).

\paragraph{Further details.} The SI documents the data and crosswalks (\SI{1}), the prestige measures (\SI{2}),
the coupling estimator, its heterogeneity and the cross-source comparison (\SI{3}), the selectivity designs and
the specification curve (\SI{4}), the UK analysis (\SI{5}), career time and the calibration of its intervals
(\SI{6}), licensing and employer sector (\SI{7}), tail outcomes (\SI{8}), the power of every null and the
hypotheses examined (\SI{9}), and earnings as a placement measure (\SI{10}).

% =====================================================================================
\section*{Declarations}
\sloppy

\paragraph{Availability of data and materials.} All inputs are public: the Wapman et al.\ faculty-hiring
data \citep{wapman2022data}; College Scorecard field-of-study and institution files \citep{scorecard}; Census
PSEO earnings and flows \citep{pseo}, releases V4.13.0 (2025Q4) and V4.14.1 (2026Q2); the Department for
Education's ``Graduate outcomes (LEO): provider level data'' (2022-23 release, Explore Education
Statistics); the ORCID-derived academic mobility edges (Zenodo, doi:10.5281/zenodo.19651302, version
20260419), built from ORCID public data \citep{orcid}, with institutions resolved through the ROR data dump
v2.8 (Zenodo record 20512981); Opportunity Insights college tables \citep{chetty2020income}; and the 2023
American Community Survey 1-year public-use microdata (U.S. Census Bureau). Provenance for every input
(URL, version, access date and licence) is recorded in \texttt{data/raw/SOURCES.md} of the repository. No
input dataset is redistributed; the derived tables are regenerated by the scripts.

\paragraph{Code availability.} The numbers in this article come from seeded scripts in the project's
public repository (scripts 01--64; mainly 32, 35, 52 and 55--64). The figures, and the numbers marked ``SI
check'' that its header lists, come from \texttt{scripts/70\_paper\_figures.py}; the SI names the source of
every other number, including two coverage checks computed outside the scripts. The repository is at
\url{https://github.com/soray42/degrees-of-separation}.

% The following three statements are to be confirmed by the author before submission.
\paragraph{Competing interests.} The author declares no competing interests.

\paragraph{Funding.} Not applicable.

\paragraph{Authors' contributions.} S.Y.\ designed the study, carried out the analysis and wrote the
manuscript.

\bibliographystyle{plainnat}
\bibliography{refs_merged}

\end{document}
